%% kapa: exam=1 round=22 year=2011 period=1 form=A subject=civil_law subject_ko=민법 seq=13 number=13
%% source: Q-Net 감정평가사 기출문제 (cst00309) artlSeq=5036790 file=extracted/1008/5036790/2011년 제22회 감정평가사 1차시험 기출문제/1교시.hwp
%% conversion: kordoc 4.24.51 -> tools/gen1.py (auto); review=proofread-sonnet; flags=none
\begin{question}{13}
乙은 甲소유의 토지를 20년간 소유의 의사로 평온ㆍ공연하게 점유하였음을 이유로 甲에게 소유권이전등기를 청구할 수 있게 되었다. 다음 설명으로 옳지 않은 것은? (다툼이 있으면 판례에 의함)
\begin{choices}
  \item 乙이 취득시효기간의 만료로 소유권이전등기청구권을 취득한 후 점유를 상실하였다면 바로 소유권이전등기청구권은 소멸한다.
  \item 乙이 등기를 경료하지 않고 있는 사이에 甲이 丙에게 그 토지를 처분하고 이전등기까지 마쳤다면, 특별한 사정이 없는 한 乙은 취득시효완성을 가지고 丙에게 대항할 수 없다.
  \item 乙이 등기를 경료하지 않고 있는 사이에 甲이 丙에게 그 토지를 처분하여 이전등기를 해 준 후 甲이 다시 토지소유권을 취득한 경우, 乙은 甲에게 시효완성을 주장할 수 있다.
  \item 乙이 등기를 경료하지 않고 있는 사이에 甲이 丙에게 그 토지를 처분한 후 乙이 소유의 의사로 평온ㆍ공연하게 다시 20년 넘게 점유한 경우, 乙은 丙에게 취득시효의 완성을 주장할 수 있다.
  \item 만약 甲이 취득시효 진행 중에 그 토지를 丙에게 양도하고 등기를 이전해 준 후에 시효가 완성되었다면, 乙은 丙에게 시효취득을 주장할 수 있다.
\end{choices}
\end{question}
\answer{1}
